\documentclass[11pt]{article}

\usepackage[a4paper,margin=1in]{geometry}
\usepackage{amsmath,amssymb,amsthm}
\usepackage{graphicx}
\usepackage{microtype}
\usepackage[hidelinks]{hyperref}

\newtheorem{theorem}{Theorem}

\title{The Full Strong Completeness of $L^p(T)$ Is Settled}
\author{Literature-answer packet for arXiv:2304.07251\\
Run \texttt{fa\_banach\_001}, lane 7}
\date{12 August 2026}

\begin{document}
\maketitle

\begin{center}
\fbox{\parbox{0.91\textwidth}{\textbf{Status.}
Already answered in the literature.  Azouzi's later arXiv:2512.11809 proves
strong completeness for nets in $L^p(T)$ for every $p\in[1,\infty]$, under
the same conditional-Riesz-triple hypotheses.}}
\end{center}

\section{The source question}

Azouzi's 2023 paper proves \emph{sequential} strong completeness of
$L^p(T)$ for every $p\in[1,\infty]$.  It explicitly distinguishes this from
full strong completeness, where arbitrary Cauchy nets rather than only
sequences must converge.  The paper records that the full problem was known
only at $p=2$ and $p=\infty$ \cite[p.~2]{Azouzi2023}.

\begin{figure}[ht]
  \centering
  \includegraphics[width=\textwidth]{figures/open_problem_crop.png}
  \caption{The open full-completeness problem in arXiv:2304.07251, PDF page 2.}
\end{figure}

The setting is a conditional Riesz triple $(E,e,T)$: $E$ is a Dedekind
complete Riesz space with weak order unit $e$, and $T$ is a conditional
expectation operator.  The vector-valued $p$-norm is
\[
  N_p(x)=T(|x|^p)^{1/p},
\]
and strong completeness means convergence of every net that is Cauchy for
this norm in the order sense.

\section{The later answer}

The same author returned to this exact problem in the 2025 preprint
\emph{Strong completeness of $L_p$-type vector lattices}
\cite{Azouzi2025}.  Its Theorem~19 proves, for each fixed
$p\in(1,\infty)$, the equivalence
\[
  L^1(T)\text{ is strongly complete}
  \quad\Longleftrightarrow\quad
  L^p(T)\text{ is strongly complete}.
\]
Combining this with the already established $p=2$ theorem yields the full
finite-$p$ range; the known $p=\infty$ result supplies the endpoint.  The
paper states the resulting conclusion as Theorem~20.

\begin{figure}[ht]
  \centering
  \includegraphics[width=\textwidth]{figures/supporting_answer_crop.png}
  \caption{Theorem~20 of arXiv:2512.11809, PDF page 16.}
\end{figure}

\begin{theorem}[Azouzi, 2025]
Let $(E,e,T)$ be a conditional Riesz triple.  Then $L^p(T)$ is strongly
complete for every $p\in[1,\infty]$.
\end{theorem}

This is precisely the property left open in arXiv:2304.07251.  It concerns
arbitrary nets, not merely sequences, and it includes all formerly missing
values of $p$.

\section{Match audit}

The source and answer align without a change of scope:
\begin{itemize}
  \item both assume a conditional Riesz triple $(E,e,T)$;
  \item both use the same spaces $L^p(T)$ and their natural vector-valued
        norms;
  \item the source asks for full strong completeness of nets;
  \item the later theorem asserts full strong completeness for every
        $p\in[1,\infty]$.
\end{itemize}

Accordingly, the open question should be marked completely answered in the
literature, rather than treated as a target for a new proof.

\section{Human review recommendation}

This is a direct theorem-to-question match and needs only light review.
Confirm the definition of ``strongly complete'' in arXiv:2512.11809 (every
strong Cauchy net converges) and the common conditional-Riesz-triple
hypotheses.  Theorem~20 then closes the source question verbatim.

\begin{thebibliography}{9}
\bibitem{Azouzi2023}
Youssef Azouzi,
\emph{Uniform Integrability in vector lattices and applications},
arXiv:2304.07251, 2023.
\url{https://arxiv.org/abs/2304.07251}.

\bibitem{Azouzi2025}
Youssef Azouzi,
\emph{Strong completeness of $L_p$-type vector lattices},
arXiv:2512.11809, 2025.
\url{https://arxiv.org/abs/2512.11809}.
\end{thebibliography}

\end{document}
